\section{方法论}
\label{sec:methodology}

\subsection{系统架构}
FruitTreeScanner系统由三个主要层组成：数据采集层、处理层和产量估算层。数据采集层利用消费级设备集成的LiDAR传感器和RGB相机，与ARKit配合实现实时点云流和空间追踪。处理层执行点云预处理、颜色过滤、聚类和多模态融合。产量估算层实现体积计算和遮挡校正以产生产量估算。

\subsection{数据采集模块}
LiDAR传感器以约60帧每秒的速度捕获三维点云数据，每帧包含约30,000个深度点。RGB相机同时记录与点云空间对齐的视觉信息。通过ARKit的传感器融合功能实现深度和颜色数据之间的时间同步。这使得带颜色的点云生成成为可能，其中每个点都携带空间坐标（x, y, z）和颜色信息（R, G, B）。

ARKit框架提供世界追踪功能，能够在帧之间维持坐标系，使得扫描移动过程中点云累积成为可能。与单帧捕获相比，这有助于构建更完整的扫描场景3D表示。

\subsection{点云预处理}
预处理操作准备原始点云数据用于聚类分析。噪声过滤去除可能是测量伪影或背景噪声的孤立点。统计离群值移除识别其邻居距离显著大于局部平均距离的点。

颜色空间转换将RGB颜色信息转换为更适合果实检测的颜色空间。系统支持HSV（色相、饱和度、值）和Lab颜色空间表示，它们将色度信息与强度分离，并提供更直观的颜色分割阈值。

降采样减少点云密度以降低计算负担，同时保留几何特征。体素网格降采样将空间划分为体素箱，并从每个箱中保留代表性点。

\subsection{基于颜色的过滤}
点云中的果实检测依赖于颜色阈值来识别可能包含果实的区域。不同果实类型表现出特征性的颜色范围，可以在HSV颜色空间中表示为色相间隔。例如，成熟苹果通常呈现红色色相值（约0-15度），而柑橘类水果显示橙色至黄色色相。

过滤算法评估每个点的颜色与目标果实类型的预定义阈值。处于指定颜色范围内的点被保留为果实候选，而落在范围外的点被分类为非果实背景。

\subsection{自适应DBSCAN聚类}
自适应DBSCAN算法通过基于局部点密度和距传感器距离动态计算邻域半径$\epsilon$来扩展标准DBSCAN。这种适应性解决了LiDAR数据的特性，即点密度随距传感器距离增加而减小。

自适应$\epsilon$计算考虑两个因素：局部点密度和深度相关缩放。距传感器更远的点需要更大的邻域半径来捕获与近点相当的局部密度。算法将$\epsilon$计算为k最近邻距离的函数，并按点的深度进行缩放。

聚类验证在聚类后进行，以过滤掉不太可能对应于果实的虚假簇。验证标准包括簇大小约束（太少点不能代表果实，太多点可能表示合并的对象或背景）、形状规则性指标（果实接近球形）和簇内颜色一致性。

\subsection{多模态融合策略}
融合策略将来自基于CoreML的图像分析的2D检测结果与来自点云处理的3D聚类候选相结合。这种方法利用了互补优势：2D检测基于视觉特征提供稳健的分类，而3D聚类提供准确的空间定位和尺寸估计。

融合过程使用与每个检测相关联的深度信息将2D检测边界框投影到3D空间中。投影区域内包含点云簇作为候选匹配。每个匹配的置信度分数考虑检测置信度、2D和3D区域之间的空间重叠以及与预期果实尺寸的尺寸一致性。

没有对应3D簇的未匹配2D检测可能表示被遮挡或部分可见的果实。没有2D检测支持的未匹配3D簇被评估为可能的检测器漏检的新型果实候选。

\subsection{遮挡校正模型}
遮挡校正解决了在任何单一扫描位置都看不到隐藏在树冠内的果实的挑战。系统采用基于可见部分估计的统计校正方法。

校正模型假设冠层内果实分布可以近似为球壳模型。可见表面比例取决于冠层密度和观察几何。通过分析果实颜色区域中可见点与总点的比例并与预期分布模式进行比较，模型估计被遮挡果实的比例。

\subsection{产量估算}
系统实施双路线产量估算策略。基于体积的路线从点云簇几何计算单个果实体积，使用果实类型特定的密度参数近似果实重量，并跨所有检测到的果实进行汇总。冠层回归路线采用统计模型将冠层体积或表面积与预期果实负载相关联，并基于可见果实计数进行校准。

最终产量估算结合两条路线的输出，根据置信度水平进行加权。对于具有完整点覆盖的充分暴露的果实，基于体积的路线获得更高的权重。对于仅有部分信息可用时被遮挡的区域，冠层回归提供补充估算。